\section{The class-sum span on the original base tuple}

% BEGIN F76B1
\paragraph{F76B1. Raw generators at the actual bad modulus and source rank.}
Fix the original second-field core $T,R,F$ with its core laws,
weight and vector tie, and a pinned base record $O$ at row $i$ and
tuple $\mu_0$. Write $K=O.c'.F\subset\mathbb C$, $j=i+1$,
$X_o=T.Q.Xo(j)$, $A=T.Q.Dw(j).\mathrm{La}$, $M=T.Q.l.M$,
$b=O.\mathrm{bad}(*)$, $r=D.R.r$ and $n=\operatorname{rank}(V_i)$.
Use the given embedding $K\hookrightarrow\mathbb C$. The rank
and weight ties give $n=D.R.nV=\kappa_i.n$, $a=\mathrm{wtA}(\kappa_i)$
dominant and $n$ even. Assume $8D.R.f\,r+4\le D.R.nV$.
Then every base cusp section over $K$ is a finite $K$-linear
combination of raw sections with generating residue tuples. These
sections have unimodular row lattice $\Lambda$, trivial at both
$b$ and $M$, allowed level $N$, and polynomial in
$\mathrm{harmL}(O.c',\Lambda,a)$. The tuple in each raw
generator may be chosen to be the same tuple of the construction
below, before the later column moves.

Here the chart argument must be made over $O.c'$. The model
$X_o$, host, level, geometric datum and integer ideals are those
of the original row. The point ring is extended from $T.Q.c.F$
by $O.\mathrm{hle}$. For its given inclusion in $\mathbb C$ the
good-embedding condition follows from the lower inclusion.
The uniformization used to define the cusp space is the original
one with its coefficient map extended. Reading a cusp section
there is bounded by the definition of $\mathrm{cuspOver}$.
Construct the auxiliary principal atlas of C76P1 with coefficient
record $O.c'$ and the indicated embedding: PCL1 supplies the same
rational Hermitian space, and the section bases of D74U2 commute
with this extension. NFL22 applies to this coefficient record,
giving the common scalar and matrix, the principal groups, and
trivial characters. C76P2 transfers the bound of the given
section from the original atlas to this atlas on the same complex
fiber. Thus it suffices to prove the relation on the auxiliary
atlas and descend it as a relation of sections. This avoids any
choice of new dictionaries for $O$.

Repeat the fixed-chart proof of B76S2 at auxiliary finite
integer $M_0=2bM$. The tuple in C76P4 is fixed before selecting
the lattice in the theta expansion. Its first $r$ entries generate
$A/IA$ for $I=\mathrm{Mfr}(l)$ and its other entries vanish.
This satisfies $\mathrm{AdmX}(\mathrm{true},x^0)$ since
the images of all $n$ entries generate; the identity cover here is
the one with threshold $n_V=n$. In particular the phase on the
identity cover is the point character of this tuple. The C76P5
level tie and the stabilizer in C76P7 are unaffected by $M_0$:
their ideals and groups are defined at $l.M$. The new finite
integer is positive and belongs to $I$. L76F6 and I75K3 therefore
apply with the majorant of F74N3 at this $M_0$. The size
inequality for them is the one displayed above because
$[E:\mathbb Q]=2D.R.f$. For each resulting $\Lambda$ its
agreement with the standard lattice at $2bM$ gives a positive
$t$ prime to $2bM$ carrying each lattice into the other by
multiplication. In particular $\mathrm{TrivAt}(b)$ and
$\mathrm{TrivAt}(M)$ hold, and $N=Mt$ is an allowed odd level
as in B76S2. The polynomial renaming, integral polynomial span,
and component constants in that proof depend only on the component
or the indicated generator in exactly the same order at $2bM$.
A76T2 gives the sections over $K$ with the prescribed short rules.
The comparison S76G4 gives the complex span of their readings on
the whole auxiliary atlas. Independence after reading (ICL8) gives
a relation over $K$ among the sections themselves.

Finally multiply each such generator by the nonzero scalar
$\beta_H(\Lambda)\in K$ as in S74B3. On the identity cover its
normalized short rule now gives precisely the raw rule, with
Fourier weight $\widehat{\mathbf1}_{\bar x^0}$: in the finite
Fourier sum the point indicator leaves its single level character.
The quotient tuple $\bar x^0$ generates $A/IA$. This verifies the
raw input used by the base cusp span, including the original test
rings and their framed points. The section relation may now be
evaluated on the original core atlas, in particular on $O.\mathrm{dict}$.
% END F76B1

% BEGIN F76B2
\paragraph{F76B2. From the raw span to class sums over the pinned base field.}
Retain the fixed data of F76B1 and the native oscillator extension
$\mathrm{IoExt}(R)$. Let $\mathcal S$ be the set of original cusp
sections $s$ over $K$ for which some coefficient vector $x$ satisfies
both $O.\mathrm{Plc}(*,x)$ and $\mathrm{IsSum}(O,*,x,s)$.
Then the original base-tuple statement is
\[
          \mathrm{cuspOver}(O,\kappa_i)
                       =\operatorname{span}_{K}\mathcal S .
\]

For the upper-level sections used here the coefficient map is
$O.j_L:\mathrm{LL}\to K$. Its composite with $K\hookrightarrow
\mathbb C$ is $\mathrm{embLL}$ by the lift pin, so each prescribed
embedded coefficient belongs to this very $K$. The construction
of T73C1--T73C5 on this fixed record uses the available upper
cover from its pins. More explicitly the tuple pins give the
half-level character and Fourier weights, representatives, lattice
bounds and normalizing scalars. The character values away from
the specified bad primes belong to $K$: the level and family
laws identify them with the corresponding embedded coefficient
values, and $O.j_L$ gives their preimages. The cover comparison
constant and the finite character have finite orders; the product
of those orders is a common exponent. These are the inputs to
the character descent in T73C2--T73C3. The model tie supplies
the dominant weight and the last positive row condition where
cuspidality needs it. T73C4 gives the cusp condition in the
prescribed extended lower uniformizations. Thus for each
placed coefficient tuple the class-sum section is available
already on $O$. This argument uses the level and family laws
from the fixed core. It does not require selecting cover or
coefficient fields for other tuples.

Let $L=\mathrm{atDict}\circ\mathrm{rdU}$ for the original record.
For each raw generator of F76B1 use the residue move of S74B3 to
the tuple $x_1$ supplied by the finite-point indicator of S74B1.
The half-level pin supplies this generating tuple. Next use the
split move of S74B4. Both moves preserve the allowed finite levels
at $b,M$ as required there; the lattice, polynomial and weakly
equivalent short level move together. V74M3 compares the old and
new normalization-one values on every $O.\mathrm{dict}(g)$ by
one nonzero scalar per move, independent of $g$. In particular
these steps need only a complex scalar, since they are being
applied after $L$. The final polynomial comes from $O.\mathrm{io}$
by the harmonic spanning pin. P74I2--P74I7 apply with the
given oscillator extension, stored extremal vector, marked target
movers and the finite tuple data of this record.

Now J74T1--J74T2 place each moved reading in the complex span
of $L(\mathcal S)$. Indeed the single-class kernel is a finite
linear combination of lifts of rational source inputs by K73G3.
Their coordinates from the pinned representatives satisfy
$O.\mathrm{Plc}$, and the preceding geometric construction gives
the corresponding $s_j\in\mathcal S$. The single-class and
integrated kernel identities on the same original dictionary carry
the same nonzero frame factor $C_{k(g)}\mu(g)$. Cancelling it gives
the moved short value as
$(c_0\sigma(\mathrm{epF}(L_h)))^{-1}\sum_j b_j Ls_j$.
The denominator is nonzero by J74T1 and the root-of-unity pin.
The cancellation does not require a nonzero value of any section
at the chosen point. F76B1 and the two moves therefore give
$Lf\in\operatorname{span}_{\mathbb C}L(\mathcal S)$ for every
$f\in\mathrm{cuspOver}(O,\kappa_i)$. K73G4 descends through
this particular dictionary: an independent finite family of
sections over $K$ has independent complex readings there.
Taking a basis from the finitely many $s_j$ gives the asserted
inclusion over $K$. The reverse inclusion follows from the
definition of $\mathcal S$ and the submodule property of cusp
sections.
% END F76B2

% BEGIN F76B3
\paragraph{F76B3. Range of the base result in the rational lift comparison.}
The base-span input of RTG16 is now supplied by F76B2 for an
original pinned base with the above core, weight tie, oscillator
extension and rank bound. If an original all-tuple record with
the comparison (RTG.17) is supplied, its class-sum existence and
the two point identities with common frame character and level
action give $\mathrm{LiftsSpan}$ exactly as in RTG16. To see
where the base enters, write a cusp form on the common geometric
target as a finite $K$-linear combination of the sections in
F76B2. The placed vector for each base section lifts to a rational
source vector by the base presentation. Multiplying its reading
by the fixed comparison character and applying the level action
expresses it as a complex scalar times the base lift of that
source vector. Applying this to the geometric form representing
a lift from the other tuple proves membership in the span of
those base lifts. RTG17--RTG18 transport the prescribed phase
only on their stated field, base-map and phase hypotheses.
F76B2 concerns a given $O$ and row; it supplies the class-sum
span before a choice of the base-map phase. Existence of the
records and ties quantified over families in the rational-tests
obligation is a separate step.
% END F76B3
